\section*{Zahvalnica}
Zahvaljujem se...

\newpage

\section*{Rezime}

Ubrzani razvoj veštačke inteligencije i velikih jezičkih modela (LLM) menja praksu provere znanja u visokom obrazovanju, naročito u ocenjivanju pitanja otvorenog tipa.
Tradicionalno ocenjivanje je vremenski zahtevno, sklono nedoslednosti i teško skalabilno za velike grupe studenata.
Veliki jezički modeli, zahvaljujući sposobnosti razumevanja i rezonovanja o prirodnom jeziku, pružaju priliku za unapređenje efikasnosti, doslednosti i učestalosti povratnih informacija.
U ovom radu predstavljen je EduGrader, višemodelska platforma otvorenog koda za ocenjivanje koja objedinjuje više velikih jezičkih modela: GPT-4o, GPT-4o-mini, DeepSeek-Chat, GPT-5.1 i DeepSeek-Reasoner.
EduGrader podržava tri nivoa strogosti ocenjivanja (blag, neutralan, strog) i generiše numeričke ocene sa sažetim obrazloženjima koja ističu ispravno rezonovanje, koncepte koji nedostaju i greške.
Sistem radi u dva režima: režimu sa zadatim referentnim odgovorom (PRA), vođenom rešenjima nastavnika, i režimu sa generisanim referentnim odgovorom (GRA), koji referentne odgovore automatski konstruiše iz nastavnih materijala.
Svi eksperimenti sprovedeni su na srpskom --- jeziku sa ograničenim resursima, slabo zastupljenom u podacima za obučavanje savremenih modela --- čime se istraživanje ocenjivanja uz podršku veštačke inteligencije proširuje izvan dominantnog engleskog konteksta.
EduGrader je evaluiran na 686 autentičnih studentskih odgovora sa šest univerzitetskih kurseva na dve institucije, uključujući tekstualna objašnjenja i odgovore u obliku koda.
U najboljoj konfiguraciji sistem postiže Pirsonov koeficijent korelacije od 0{,}90 sa ljudskim ocenjivačima, uz smanjenje opterećenja nastavnika i poboljšanje pouzdanosti, pri čemu je namenjen kao pomoć nastavnicima, a ne kao njihova zamena.

\medskip
\noindent\textbf{Ključne reči:} automatsko ocenjivanje, obrazovanje, veliki jezički modeli, veštačka inteligencija.
\newpage
